\documentclass[12pt]{article}
\usepackage[margin=0.75in,top=0.8in,bottom=0.65in]{geometry}
\usepackage{lmodern}
\usepackage{microtype}
\usepackage{fancyhdr}
\usepackage{titlesec}
\usepackage{enumitem}
\usepackage{lastpage}
\usepackage[hidelinks]{hyperref}

\setlength{\parindent}{0pt}
\setlength{\parskip}{0.38em}
\setlength{\headheight}{26pt}
\titleformat{\section}{\large\bfseries\scshape}{}{0pt}{}
\titleformat{\subsection}{\normalsize\bfseries}{}{0pt}{}
\pagestyle{fancy}
\fancyhf{}
\fancyhead[C]{\footnotesize Lit Example Reader \quad Assignment ID: U1L6LE \quad Page \thepage\ of \pageref{LastPage}}
\renewcommand{\headrulewidth}{0.5pt}

\newenvironment{playpassage}{\begin{quote}\small\setlength{\parskip}{0.25em}}{\end{quote}}
\newcommand{\speaker}[1]{\textbf{#1.}}

\begin{document}

\begin{center}
{\LARGE\bfseries Week 6 Lit Example Reader}\par\vspace{0.05em}
{\large\scshape Julius Caesar: Valid Form, Unproved Premise}\par\vspace{0.15em}
\rule{0.78\textwidth}{0.5pt}
\end{center}

\textbf{Purpose:} This reader gives Julius Caesar passages for applying Week 6's logic habit: separate validity from truth and soundness. A valid argument can still fail if an important premise is false, doubtful, or unproved.

\textbf{What to watch for:} Brutus can arrange his defense into a valid shape. Antony's reply matters because it challenges the factual pressure point. Week 6 asks you to run both tests: form first, fact second.

\section*{Before the Passages: The Week 6 Logic Tool}

\begin{itemize}[leftmargin=1.5em,itemsep=0.2em]
\item \textbf{Form test:} If the premises were true, would the conclusion have to follow?
\item \textbf{Fact test:} Are the premises actually true, adequately supported, or at least reasonable?
\item \textbf{Sound argument:} valid form plus true premises.
\end{itemize}

Do not call an argument sound simply because the conclusion sounds important. Do not call it invalid simply because a premise is doubtful. Keep the two tests separate.

\section*{Passage 1: Act 3, Scene 2 --- Brutus's Valid-Looking Defense}

\textbf{Context:} Brutus tells Rome why Caesar was killed. The speech can be reconstructed in a valid form, but the form does not prove every premise.

\begin{playpassage}
\speaker{Brutus} If then that friend demand why Brutus rose against Caesar,\par
this is my answer:---Not that I loved Caesar less, but that I\par
loved Rome more. Had you rather Caesar were living and die all\par
slaves, than that Caesar were dead, to live all free men?\par
As Caesar loved me, I weep for him; as he was fortunate, I rejoice\par
at it; as he was valiant, I honour him: but, as he was ambitious,\par
I slew him. There is tears for his love; joy for his fortune;\par
honour for his valour; and death for his ambition.\par
Who is here so base that would be a bondman? If any, speak;\par
for him have I offended. Who is here so rude that would not be a Roman?\par
If any, speak; for him have I offended. Who is here so vile that\par
will not love his country? If any, speak; for him have I offended.
\end{playpassage}

\subsection*{Reader Notes}

Here is your most assisted example. Brutus's argument can pass a form test if reconstructed this way: All men whose ambition would enslave Rome deserve death for Rome's freedom; Caesar was a man whose ambition would enslave Rome; therefore Caesar deserved death for Rome's freedom. If both premises are true, the conclusion follows. But the minor premise still needs the fact test.

\textbf{Reading task:} Mark Brutus's conclusion. Box the premise about Caesar's ambition. In the margin, write: form test or fact test?

\newpage

\section*{Passage 2: Act 3, Scene 2 --- Antony's Fact Test}

\textbf{Context:} Antony does not mainly say Brutus's form is invalid. He gives reasons to doubt whether Caesar fits Brutus's pressure-point premise.

\begin{playpassage}
\speaker{Antony} The noble Brutus\par
Hath told you Caesar was ambitious:\par
If it were so, it was a grievous fault,\par
And grievously hath Caesar answer'd it.\par
He was my friend, faithful and just to me:\par
But Brutus says he was ambitious;\par
And Brutus is an honourable man.\par
He hath brought many captives home to Rome\par
Whose ransoms did the general coffers fill:\par
Did this in Caesar seem ambitious?\par
When that the poor have cried, Caesar hath wept:\par
Ambition should be made of sterner stuff:\par
Yet Brutus says he was ambitious;\par
And Brutus is an honourable man.\par
You all did see that on the Lupercal\par
I thrice presented him a kingly crown,\par
Which he did thrice refuse: was this ambition?\par
Yet Brutus says he was ambitious.
\end{playpassage}

\subsection*{Reader Notes}

Antony's evidence does not automatically prove the opposite of Brutus's claim. It may make the claim doubtful. That distinction matters. A premise can fail the fact test because it is false, but it can also fail because the speaker has not given enough evidence for it.

\textbf{Reading task:} Number Antony's evidence. For each item, write whether it proves Caesar was not ambitious or merely gives a reason to doubt Brutus's charge.

\section*{How to Use These Passages on the Worksheet}

The worksheet asks for a final diagnosis: valid, invalid, sound, unsound, or not yet proved. Use those labels carefully. A strong answer can say, ``valid form, but not sound unless the ambition premise is true.''

\end{document}
